@@M0@@. Let @@M1@@ be the connected components of @@M2@@; then @@M3@@ decomposes, @@M4@@, and the assertion follows. Finally, an immediate induction reduces us to the case of a morphism @@M5@@. After shrinking @@M8@@, @@M6@@ and @@M7@@ then follow, by (EGA III 7.8.6), from the

\sgaStableTarget{sga6.e12.prop.2-3}\textbf{Proposition 2.3.} \uline{Let} @@M9@@ \uline{be a scheme,} @@M10@@ \uline{a proper morphism, and} @@M11@@ \uline{two closed subschemes defined by the ideals} @@M12@@\uline{. Suppose that}
\begin{enumerate}[label=\alph*)]
\item @@M13@@ \uline{are faithfully flat and of finite presentation over} @@M14@@.
\item @@M15@@ and @@M16@@ \uline{universally (where} @@M17@@, @@M18@@).
\item @@M19@@.
\end{enumerate}
\uline{Then} @@M20@@ \uline{is representable by an affine morphism of finite presentation.}

Indeed, first note that the hypotheses are stable under every base change. For (c), this follows from the two exact sequences:
@@M21@@
and
@@M22@@
the first shows that @@M23@@ is flat over @@M24@@; then the second shows that @@M25@@ is stable under base change.

The assertion is local for the f.p.p.f. topology. Making the base change @@M26@@, we may therefore suppose that @@M27@@ has a section @@M28@@, which defines sections @@M29@@ of @@M30@@ and @@M31@@ of @@M32@@.

After a suitable base change, we are reduced to seeing that every point @@M33@@ has as its inverse image a subfunctor of

